\section{Methodology}

\subsection{Dataset Description}

The dataset was provided by the Maersk Mc-Kinney Moller Institute at the University of Southern
Denmark as part of the Tools of Artificial Intelligence course. It contains single-channel EMG
recordings for three hand gestures (Rock, Paper, Scissor) stored as separate CSV files. For this
binary classification task, \texttt{RockData.csv} (2\,909 trials, label\,=\,0) and
\texttt{PaperData.csv} (2\,941 trials, label\,=\,1) were merged to form a combined dataset of
5\,850 trials. Each trial contains 1\,024 time-series samples recorded at a sampling frequency
of $f_s = 256$\,Hz, corresponding to approximately 4\,seconds of signal per trial. The class
distribution is nearly balanced (49.7\% Rock, 50.3\% Paper), eliminating the need for any class
balancing technique. Figure~\ref{fig:visualisation} shows the exploratory visualisation of the
dataset.

\begin{figure}[H]
    \centering
    \includegraphics[width=0.82\textwidth]{figures/visualisation}
    \caption{Exploratory dataset visualisation. (Top-left) Representative EMG signals for Rock
    and Paper trials. (Top-right) Class distribution showing near-balanced classes (Rock: 2\,909,
    Paper: 2\,941). (Bottom-left) Mean signal per class averaged over all trials. (Bottom-right)
    Amplitude density distributions per class, showing significant overlap reflecting
    single-channel measurement limitations.}
    \label{fig:visualisation}
\end{figure}

\subsection{Data Preprocessing}

\textbf{Missing Value Analysis.}
The complete feature matrix ($5\,850 \times 1\,024$) was examined for missing values. The total
NaN count was zero, confirming a complete dataset requiring no imputation. Had missing values
been present, mean imputation would have been applied --- it preserves the signal's overall
amplitude distribution and is appropriate when missingness is small and approximately random
\citep{Little2002}.

\textbf{Signal Filtering.}
According to the dataset description provided by the course, the EMG signals were already
filtered prior to recording. Therefore, no additional signal processing (notch filtering at
50\,Hz or bandpass filtering between 0.5--30\,Hz) was applied. The signals are used directly
for feature extraction.

\subsection{Feature Extraction}

A total of 26 features were extracted from each trial across four domains. Multi-domain feature
extraction is motivated by the complementary nature of different feature types: statistical
features capture amplitude distribution, time-domain features capture temporal dynamics,
frequency-domain features capture spectral content, and entropy features capture signal
complexity. All features were computed on the full 1\,024-sample trial signal.

\begin{itemize}
    \item \textbf{Statistical features (8):} Mean, Median, Standard Deviation, Variance,
          Skewness, Kurtosis, Range, Interquartile Range (IQR).
    \item \textbf{Time-domain features (7):} Root Mean Square (RMS), Zero Crossing Rate (ZCR),
          Autocorrelation, Mean Absolute Deviation (MAD), Maximum, Minimum, Signal Energy.
    \item \textbf{Frequency-domain features (6):} Dominant Frequency, Total Power,
          Band Power (0.5--40\,Hz), Mean Frequency, Median Frequency, Frequency Variance ---
          computed using Welch's Power Spectral Density estimate at $f_s = 256$\,Hz.
    \item \textbf{Entropy features (5):} Sample Entropy, Spectral Entropy, Permutation Entropy,
          SVD Entropy, and Lempel-Ziv (LZiv) Complexity --- quantifying signal regularity and
          complexity.
\end{itemize}

Key feature definitions are provided below. The mean and standard deviation characterise signal
amplitude distribution:

\begin{equation}
    \mu = \frac{1}{N} \sum_{i=1}^{N} x_i
    \label{eq:mean}
\end{equation}

\begin{equation}
    \sigma = \sqrt{\frac{1}{N} \sum_{i=1}^{N} \left(x_i - \mu\right)^2}
    \label{eq:std}
\end{equation}

\noindent where $x_i$ is the $i$-th sample and $N = 1\,024$ is the number of samples per trial.
RMS and Signal Energy quantify signal power:

\begin{equation}
    \mathrm{RMS} = \sqrt{\frac{1}{N} \sum_{i=1}^{N} x_i^2}
    \label{eq:rms}
\end{equation}

\begin{equation}
    E = \sum_{i=1}^{N} x_i^2
    \label{eq:energy}
\end{equation}

\subsection{Feature Selection}

\textbf{Step 1 --- Standardisation.}
All 26 features were standardised to zero mean and unit variance using \texttt{StandardScaler}
prior to feature selection, as both Kruskal-Wallis and PCA are sensitive to feature scale
differences.

\textbf{Step 2 --- Kruskal-Wallis Test.}
The non-parametric Kruskal-Wallis H-test was applied to each feature to assess whether its
distribution differed significantly between Rock and Paper classes. Features with $p < 0.05$
were retained. The Kruskal-Wallis test was chosen over ANOVA because it does not assume
normality --- a critical consideration for entropy-based measures. Figure~\ref{fig:kruskal}
shows the four most discriminative features identified.

\begin{figure}[H]
    \centering
    \includegraphics[width=0.78\textwidth]{figures/kruskal_boxplot}
    \caption{Box plots of the four most discriminative features identified by the Kruskal-Wallis
    test ($p < 0.05$). Rock class (blue) shows consistently higher median values than Paper
    class (orange) for all four features, indicating clear distributional differences that
    support classification.}
    \label{fig:kruskal}
\end{figure}

\textbf{Step 3 --- PCA.}
Principal Component Analysis was applied to the selected features, retaining components that
cumulatively explain 95\% of the total variance. PCA transforms the feature space into
orthogonal, uncorrelated components ordered by variance, removing redundant information and
reducing the risk of overfitting. Figure~\ref{fig:pca} shows the 2D PCA projection.

\begin{figure}[H]
    \centering
    \includegraphics[width=0.52\textwidth]{figures/pca_scatter}
    \caption{2D PCA projection of the gesture feature space after Kruskal-Wallis selection.
    Rock trials (blue) and Paper trials (orange) show partial but meaningful separation along
    PC\,1. The overlap reflects the inherent difficulty of single-channel EMG discrimination
    between these two gestures.}
    \label{fig:pca}
\end{figure}

\subsection{Classification Models}

Three classifiers representing distinct algorithmic families were implemented and compared:

\begin{itemize}
    \item \textbf{Random Forest (RF):} An ensemble of 100 decision trees trained on bootstrap
          samples. Each tree independently votes and the majority determines the prediction. RF
          is robust to overfitting and handles non-linear boundaries well \citep{Breiman2001}.
    \item \textbf{SVM with RBF Kernel (SVM-RBF):} Maps the feature space to a
          higher-dimensional space using the kernel
          $K(\mathbf{x}_i, \mathbf{x}_j) = \exp\!\left(-\|\mathbf{x}_i - \mathbf{x}_j\|^2 /
          2\sigma^2\right)$, enabling non-linear classification while maximising the margin
          between classes \citep{Oskoei2007}.
    \item \textbf{XGBoost:} A regularised gradient-boosted tree ensemble that iteratively fits
          trees to the residual errors of previous trees. Includes L1 and L2 regularisation to
          control overfitting \citep{Chen2016}.
\end{itemize}

\subsection{Validation Strategies}

\begin{itemize}
    \item \textbf{Hold-out (80/20 split):} 4\,680 training samples and 1\,170 test samples
          using stratified random sampling to preserve class proportions. Fast and simple, but
          sensitive to the random split chosen.
    \item \textbf{5-Fold Stratified Cross-Validation:} The dataset was partitioned into 5 folds
          of equal size. Each fold served once as the test set while the remaining 4 folds were
          used for training. Results are averaged across all 5 folds, providing a reliable
          estimate of generalisation.
    \item \textbf{GridSearchCV:} Exhaustive hyperparameter search applied to all three
          classifiers using 5-fold CV: RF over $n\_\mathrm{estimators}$,
          $\mathrm{max\_depth}$, $\mathrm{min\_samples\_split}$ (18 configs, 90 fits);
          SVM-RBF over $C$ and $\gamma$ (12 configs, 60 fits); XGBoost over
          $n\_\mathrm{estimators}$, $\mathrm{max\_depth}$, $\mathrm{learning\_rate}$
          (18 configs, 90 fits).
\end{itemize}
